% !Mode:: "TeX:UTF-8"
% !TEX program = xelatex
% =====================================================================
%  “华为杯”第二十三届中国研究生数学建模竞赛  参赛论文
%  2026
%
%  编译:  ./build.sh        （产物在 build/main.pdf）
%  预览:  ./preview.sh      （渲染前几页成 PNG 检查排版）
%
%  ⚠️  除封面外，全文任何地方不得出现学校 / 姓名 / 队伍编号。
%      定稿前务必自查：PDF 元数据、摘要页之后的页面不含身份信息。
% =====================================================================

\documentclass[bwprint,fontset=fandol]{gmcmthesis}
%  fontset 说明：
%    fandol  —— Linux/TeXLive 自带的开源中文字体，本机已验证可编译（推荐）
%    windows —— 需要 SimSun/SimHei/KaiTi，只有装了 Office 字体的 Windows 能用
%    mac     —— macOS

\usepackage[framemethod=TikZ]{mdframed}
\usepackage{subfig}
\usepackage{colortbl}

% 算法环境
\usepackage[noend]{algpseudocode}
\usepackage{algorithmicx,algorithm}
\floatname{algorithm}{算法}
\renewcommand{\algorithmicrequire}{\textbf{输入:}}
\renewcommand{\algorithmicensure}{\textbf{输出:}}

% 便捷宏
\newcommand{\red}[1]{\textcolor{red}{#1}}
\newcommand{\todo}[1]{\textcolor{red}{【待办：#1】}}

% =====================================================================
%  封面信息  ——  这些内容只会出现在第 1 页封面上
% =====================================================================
\title{论文题目（三号黑体，居中）}
\baominghao{26XXXXXXXXX}   % 参赛队号（11 位，以平台公布为准）
\schoolname{XX大学}        % 学校名称
\membera{队员一}           % 队员 A
\memberb{队员二}           % 队员 B
\memberc{队员三}           % 队员 C

\begin{document}

% ---------- 第 1 页：封面（不可删除，4 个 logo 不可替换） ----------
\maketitle

% ---------- 第 2 页起：摘要页（题目 + 摘要 + 关键词，≤2 页） ----------
% !TeX root = main.tex
% ============================================================
%  摘要页  ——  评委的第一印象，权重远超它的篇幅
% ============================================================
%  通知四(一)4 明文要求，摘要必须包含：
%     建模思路 / 主要方法 / 模型 / 结果与结论 / 创新点
%  篇幅 ≤ 2 页，无需译成英文。
%
%  写作要点：
%   · 一定要有具体数字。
%     「求得最优方案总成本 3.27 万元，较基准方案下降 18.4%」
%     远胜「得到了较好的结果」。
%   · 一问一段，结构与正文对应。
%   · 最后写，写三遍。
% ============================================================

\begin{abstract}

% ---- 开篇：背景 + 本文做了什么（2~3 句） ----
针对……问题，本文……。

% ---- 问题一 ----
\textbf{针对问题一}，建立了……模型。首先……，其次……，采用……求解，
得到……。结果表明……（\textbf{给出具体数字}）。

% ---- 问题二 ----
\textbf{针对问题二}，……。

% ---- 问题三 ----
\textbf{针对问题三}，……。

% ---- 检验与灵敏度 ----
最后对模型进行了……检验与灵敏度分析，结果显示……，说明模型具有较好的稳健性。

% ---- 创新点（务必单列，评委会找） ----
\textbf{本文的创新点在于}：（1）……；（2）……；（3）……。

\keywords{关键词一\quad 关键词二\quad 关键词三\quad 关键词四}

\end{abstract}


\pagestyle{plain}

% 目录：官方模板没有目录页，是否加自行决定
%\maketoc

% ---------- 正文 ----------
% !TeX root = ../main.tex
\section{问题重述}

\subsection{问题背景}

% 用自己的话写背景，不要照抄题目。点出问题的实际意义和数学本质。
% 1~2 段即可。

\subsection{问题提出}

本文需要解决以下问题：

\begin{enumerate}
  \item \textbf{问题一}：……
  \item \textbf{问题二}：……
  \item \textbf{问题三}：……
\end{enumerate}

% !TeX root = ../main.tex
\section{问题分析}

% 每一问单独一小节：这一问的难点在哪、为什么选这类模型、整体求解路线。
% 这一章配一张总体技术路线图效果最好。

\subsection{问题一的分析}

\subsection{问题二的分析}

\subsection{问题三的分析}

\subsection{总体技术路线}

% \begin{figure}[htbp]
%   \centering
%   \includegraphics[width=0.85\textwidth]{figures/fig_roadmap.pdf}
%   \caption{总体建模技术路线}
%   \label{fig:roadmap}
% \end{figure}

% !TeX root = ../main.tex
\section{模型假设}

% 每条假设都要：(1) 说清楚假设了什么 (2) 给出依据 (3) 后面真的用上
% 只为了能算而硬塞的假设会被扣分。

\begin{enumerate}
  \item \textbf{假设一}：……。
        依据：……。
  \item \textbf{假设二}：……。
        依据：……。
  \item \textbf{假设三}：……。
        依据：……。
\end{enumerate}

% !TeX root = ../main.tex
\section{符号说明}

% 只列全文反复出现的符号；一次性的局部符号在正文里就地说明即可。

\begin{table}[htbp]
  \centering
  \caption{符号说明}
  \label{tab:symbols}
  \begin{tabular}{clc}
    \toprule
    符号 & 含义 & 单位 \\
    \midrule
    $x_i$    & …… & —   \\
    $y_{ij}$ & …… & —   \\
    $\lambda$ & …… & —  \\
    \bottomrule
  \end{tabular}
\end{table}

% !TeX root = ../main.tex
\section{问题一的模型建立与求解}

\subsection{数据预处理}

\subsection{模型的建立}

% 公式用 equation 环境并编号，正文要引用它，如 式\eqref{eq:obj}。
% \begin{equation}\label{eq:obj}
%   \min_{x} \; \sum_{i=1}^{n} c_i x_i
% \end{equation}

\subsection{模型的求解}

% 算法步骤 / 伪代码。复杂度和实际运行时间要给出来。

\subsection{结果与分析}

% 给具体数字，并解释这些数字意味着什么。只贴表不解读等于没写。


\section{问题二的模型建立与求解}

\subsection{模型的建立}

\subsection{模型的求解}

\subsection{结果与分析}


\section{问题三的模型建立与求解}

\subsection{模型的建立}

\subsection{模型的求解}

\subsection{结果与分析}

% !TeX root = ../main.tex
\section{模型的检验与灵敏度分析}

% 很多队伍没时间做这一章 —— 做了就是优势。哪怕只做一问也要做。

\subsection{模型的合理性检验}

% 量纲是否正确？极端情形下模型退化成什么？结果是否与常识/文献一致？

\subsection{灵敏度分析}

% 关键参数扰动 ±10%、±20% 时结果如何变化。配图最直观。

\subsection{误差分析}

% 误差来源：数据误差、模型简化、数值求解精度。定量估计更好。

\subsection{与其他方法的对比}

% 若时间允许：与基准方法/简化模型对比，说明本文模型的增益。

% !TeX root = ../main.tex
\section{模型的评价与推广}

\subsection{模型的优点}

% 要具体。"模型科学合理" 这种话等于没写。
\begin{enumerate}
  \item ……
  \item ……
\end{enumerate}

\subsection{模型的不足}

% 诚实写不足反而加分。审稿人最烦声称自己没有缺点的论文。
\begin{enumerate}
  \item ……
  \item ……
\end{enumerate}

\subsection{模型的改进方向}

\begin{enumerate}
  \item ……
\end{enumerate}

\subsection{模型的推广}

% 这个模型还能用在什么场景。


% ---------- 参考文献 ----------
%  附件2 规定的是「中文式」格式，不是 GB/T 7714，也不是 IEEE：
%     书籍  [编号] 作者，书名，出版地：出版社，起止页码，出版年。
%     期刊  [编号] 作者，论文名，杂志名，卷期号：起止页码，出版年。
%     网络  [编号] 作者，资源标题，网址，访问时间（年月日）。
%  按正文中的引用次序排列；正文用 [1][3] 方括号标注；书籍必须给页码。
%
%  ⚠️ 自带的 gmcm.bst 输出的不是这个格式。竞赛只有 4 天，
%     最省事也最保险的做法是直接手写下面的 thebibliography 环境，
%     按正文首次引用的顺序编号。

\begin{thebibliography}{99}
\bibitem{ref1} 作者，书名，出版地：出版社，起止页码，出版年。
\bibitem{ref2} 作者，论文名，杂志名，卷期号：起止页码，出版年。
\bibitem{ref3} 作者，资源标题，\url{https://网址}，访问时间（2026年9月23日）。  % 网址用 \url{} 才能断行
\end{thebibliography}

% ---------- 附录 ----------
% !TeX root = ../main.tex
\begin{appendices}
\renewcommand{\thesubsection}{\Alph{section}.\arabic{subsection}}  % 附录小节编为 A.1、A.2

\section{核心程序代码}

% 只贴核心代码，不要把整个工程贴进来（会把篇幅撑爆，也没人看）。
% 完整程序放进上传的附件压缩包。
%
% 注意：代码里不能出现学校名、姓名、队伍编号、含身份信息的文件路径。
% 引用的开源实现 / AI 生成片段要在注释里标明来源。

\subsection{问题一求解程序}

% \begin{lstlisting}[language=Python, caption={q1\_solve.py}]
% import numpy as np
% ...
% \end{lstlisting}

\subsection{问题二求解程序}

\end{appendices}


\end{document}
